% explainer-source-hash: f749393b7f38e0d6d15c3f1e7bc1a051ad8c17439c55d82b1a6fe2bb4aefe790
% explainer-source-commit: a49871e563da2cb89c2448fb6db326242aa5301c
% explainer-generated: 2026-07-16
% Regenerate with /explain-all (see WIRING.md). Do not edit this stamp by hand:
% it is how the toolchain knows whether this document still matches the code.
\documentclass[11pt,a4paper]{article}
\input{preamble}

\title{\vspace{-1.5cm}\textbf{Paws Rescue Donations}\\[0.3cm]
\large Brief: What It Is, How It Works, and Where It Breaks}
\author{Explainer for \texttt{paws-rescue-donations}}
\date{}

\begin{document}
\maketitle
\thispagestyle{empty}

\section*{The short version}

\begin{keyidea}{In one paragraph}
A Django web application that replaces a pet rescue's donation spreadsheet with a
database behind a login. Staff record who gave how much for which animal, filter and sort
those records, watch running totals, and export the result as CSV. Built for a
faculty-supervised Khidmat community-service project at Habib University. It is small,
conventional and competently made. It also carries two real bugs, both verified by
running the code, both sitting in the only two features nobody wrote a test for.
\end{keyidea}

This document assumes no prior knowledge; a glossary at the end defines every technical
term it uses. Its companion, the Deep Dive, walks the same project function by function.

\section{Why it exists}

Small rescues track donations in a spreadsheet. That works until two or three people are
entering records, at which point three failures arrive at once: no access control, no
consistent fields, and no reliable way to ask the data a question such as ``how much has
Bella received this month''.

This project answers that with the smallest thing that could work: one database, one
login, one dashboard. There is no public side at all. No donor-facing page, no sign-up
form. Every route is behind staff authentication.

\begin{keyidea}{The complete feature set}
Email and password login; create, edit and delete donations; filter by seven criteria at
once; sort and paginate; a dashboard with running totals and a top-5-pets panel; CSV
export of the filtered view; add and remove staff logins from the web UI.
\end{keyidea}

\section{How it is built}

\begin{figure}[H]
\centering
\resizebox{0.92\textwidth}{!}{
\begin{tikzpicture}[node distance=8mm]
  \node[extbox] (b) {Browser\\ \scriptsize Bootstrap 5 via CDN};
  \node[box, below=10mm of b] (mw) {Middleware\\ \scriptsize security, whitenoise, session, CSRF, auth};
  \node[box, below=9mm of mw] (u) {URL routing};
  \node[box, below left=11mm and 12mm of u] (acc) {\code{apps.accounts}\\ \scriptsize login, staff management};
  \node[box, below right=11mm and 12mm of u] (don) {\code{apps.donations}\\ \scriptsize dashboard, CRUD, CSV};
  \node[box, below=24mm of u] (orm) {Django ORM};
  \node[store, below=9mm of orm] (db) {PostgreSQL in production\\ SQLite locally};

  \draw[flow] (b) -- node[lbl,right]{request} (mw);
  \draw[flow] (mw) -- (u);
  \draw[flow] (u) -- (acc);
  \draw[flow] (u) -- (don);
  \draw[flow] (acc) -- (orm);
  \draw[flow] (don) -- (orm);
  \draw[flow] (orm) -- node[lbl,right]{SQL} (db);
  \draw[dflow] (mw.west) to[bend right=50] node[lbl,left]{HTML} (b.west);
\end{tikzpicture}}
\caption{Everything is server-rendered. The browser receives finished HTML and runs no
JavaScript the project wrote itself.}
\end{figure}

Two Django apps carry all the behaviour. \code{accounts} decides who may log in;
\code{donations} is the subject matter. A third folder, \code{paws\_rescue}, holds
configuration rather than features. Six dependencies, all pinned: Django, psycopg2,
python-dotenv, gunicorn, whitenoise, dj-database-url.

\subsection{The whole data model}

\begin{figure}[H]
\centering
\resizebox{0.9\textwidth}{!}{
\begin{tikzpicture}
  \node[box, font=\scriptsize] (u) {{\small\textbf{User}}\\ UUID id, unique email,\\ is\_staff, is\_superuser, is\_active};
  \node[box, font=\scriptsize, right=40mm of u] (d) {{\small\textbf{Donation}}\\ integer id, donor name/phone/email,\\ pet\_name, amount (Decimal), currency,\\ payment\_method, reference\_no,\\ donation\_date, notes};
  \draw[dflow] (u) -- node[lbl,above]{this relationship\\ does not exist} (d);
\end{tikzpicture}}
\caption{Two tables, and the missing link between them: nothing records which staff
member entered a donation.}
\end{figure}

\subsection{Three decisions worth understanding}

\textbf{A custom user model, defined on day one.} Django's built-in user has a
\code{username} field; this app wants email-only login. So it defines its own user with
\code{USERNAME\_FIELD = 'email'}, a UUID primary key, and a hand-written manager to
replace the default one (whose \code{create\_user(username, ...)} signature no longer
fits). This matters more than it looks: \code{AUTH\_USER\_MODEL} is close to unchangeable
after the first migration, so it is the highest-leverage decision in a new Django
project, and getting it right at the start is a genuine credit to the build.

\textbf{Money is a \code{Decimal}, never a float.} Floats store values in binary, where
0.1 has no exact representation and \code{0.1 + 0.2} yields \code{0.30000000000000004}.
Summing thousands of donations that way accumulates real error. \code{DecimalField}
stores the digits exactly.

\textbf{Filtering rides on QuerySet laziness.} \code{dashboard\_view} stacks up to seven
filters, each applied only if the user supplied that field. None of them touch the
database. A Django QuerySet records conditions and runs nothing until something reads it,
so all seven collapse into a single SQL query with one \code{WHERE} clause.

\subsection{The dashboard, which is where the work happens}

\begin{figure}[H]
\centering
\resizebox{0.95\textwidth}{!}{
\begin{tikzpicture}[node distance=5mm]
  \node[box] (all) {\code{Donation.objects.all()}};
  \node[box, below=6mm of all] (filt) {apply up to 7 optional filters\\ \scriptsize search, date range, method, currency, amount range};
  \node[box, below=5mm of filt] (sort) {sort};
  \node[box, right=26mm of filt] (tot) {totals + top 5 pets\\ \scriptsize aggregated over the \emph{filtered} set};
  \node[box, below=6mm of tot] (pag) {paginate, 20 per page};
  \node[store, below=6mm of pag] (s) {session:\\ \scriptsize every filtered id};
  \node[box, left=22mm of s] (csv) {\code{export\_csv\_view}\\ \scriptsize reads them back};

  \draw[flow] (all) -- (filt);
  \draw[flow] (filt) -- (sort);
  \draw[flow] (sort.east) to[bend right=25] (tot.west);
  \draw[flow] (tot) -- (pag);
  \draw[flow] (pag) -- (s);
  \draw[flow] (s) -- node[lbl,above]{later GET} (csv);
\end{tikzpicture}}
\caption{One view does filtering, aggregation, pagination and the session handoff that
feeds CSV export. That last arrow is where both of the project's serious bugs live.}
\end{figure}

The CSV export link is a plain GET carrying no filter state, so the export view cannot
know what the dashboard was showing. The workaround: the dashboard writes every matching
donation id into the session, and the export view reads them back. The README already
calls this ``the weakest part of the design''. It is right, and running it shows the
weakness is more concrete than the README realised.

\section{Verification: what was actually run}

Nothing in this document is taken on trust from the README. A clean virtual environment
was built, the pinned dependencies installed, and the code executed.

\begin{lstlisting}[caption={The test suite, run on this machine}]
$ python manage.py test
Ran 11 tests in 12.828s

OK
\end{lstlisting}

Eleven tests, all passing, exactly as the README claims. They cover authentication
redirects, admin creation, the self-deletion guard, donation CRUD, search, and the CSV
response headers. The dependency pins install cleanly.

That is a real result and it is also the setup for the section below: the suite is green
and the two worst bugs are untouched by it.

\section{Where it breaks}

\begin{honest}{A filter matching nothing exports the entire database}
\code{export\_csv\_view} reads the session ids and branches on
\code{if filtered\_ids:}. An empty list is falsy, so it cannot distinguish ``the user
never loaded the dashboard, export everything'' from ``the user's filter matched zero
rows, export nothing''. Driven through the real test client with two donations present:
\begin{lstlisting}
GET /dashboard/?search=NoSuchDonorZZZ   -> page shows "0 donation(s)", session ids == []
GET /donations/export/                  -> CSV contains 2 data rows (every donor)
\end{lstlisting}
The user filters to nothing, sees ``No donations found'', clicks Export CSV, and silently
downloads every donor name, phone number and email in the system. It is triggered by an
ordinary action, it fails in the unsafe direction, and it is invisible because the file
downloads without complaint. This is the most serious defect in the project and the fix
is one line.
\end{honest}

\begin{honest}{Mixed currencies are added together and labelled PKR}
\code{Sum('amount')} ignores the \code{currency} column next to it, and the dashboard
template appends the literal text ``PKR'' to the result. With one 100 PKR donation and
one 100 USD donation present:
\begin{lstlisting}
total_all_time: 200        rendered page contains "200.00 PKR"? True
\end{lstlisting}
The dashboard reports 200.00 PKR. The rescue received 100 PKR and 100 USD. A confidently
presented wrong number is worse than an error.
\end{honest}

\begin{honest}{The README's own account of the currency issue is backwards}
The README says a free-text currency field means typos such as ``pkr'' or ``Rs'' ``would
silently split the totals''. Running it shows the opposite. Nothing splits; everything
merges, because the totals never group by currency in the first place. And
\code{currency\_\_iexact} means the filter already folds ``pkr'' into ``PKR'' correctly,
so the case-typo case is the one that works. This is worth flagging plainly: an
interviewer will read the README as the author's own assessment, and defending a
self-assessment the code contradicts is harder than having none. The real fix is not just
a choice list; even with a clean enum, summing across currencies is still meaningless.
The dashboard must group by currency or convert at a recorded rate.
\end{honest}

\begin{honest}{The development database is committed, hash and all}
\file{db.sqlite3} is tracked in git and \file{.gitignore} does not exclude it. It was
added deliberately (commit \code{a49871e}). Reading it directly: one user row for
\code{admin@example.com}, superuser, with a real PBKDF2 password hash; two expired
session rows; zero donations. The exposure is genuinely limited (placeholder email,
no donor data, expired sessions, 600{,}000-iteration hash) so this is not an emergency.
It is still a committed credential hash and a bad habit, and the next person who commits
that file after entering real donor data publishes real donor data.
\end{honest}

\begin{honest}{Every web-created admin is an unvalidated superuser}
\file{settings.py} configures the four standard password validators, but they only run
where something calls them, and \code{AdminCreationForm} never does: it subclasses
\code{ModelForm}, not Django's \code{UserCreationForm}. Verified by running the form:
a one-character password \code{'1'} validates cleanly, no errors. Meanwhile the form
hard-codes \code{is\_staff = True} and \code{is\_superuser = True} on every account it
creates, so there is exactly one privilege level. The two compound into the weakest
possible password on the most privileged possible account. The validators do run for
\code{manage.py createsuperuser}, which is why the gap is easy to miss: the settings look
correct.
\end{honest}

\begin{honest}{No audit trail}
\code{Donation} has no foreign key to \code{User}. Nothing records who entered or deleted
a record, and deletes are hard deletes with no history. For a system whose entire purpose
is tracking other people's money, this is the largest structural gap, and it is the first
thing a reviewer will ask about. The README raises it too.
\end{honest}

\begin{honest}{Smaller, and real}
\code{SECRET\_KEY} falls back to a hard-coded value published in the repo, so a
deployment that forgets the variable runs happily on a known key. The session export
grows without bound and is rewritten to the database on every request
(\code{SESSION\_SAVE\_EVERY\_REQUEST = True}). \code{donation\_date} and
\code{pet\_name} drive the ordering and the grouping and neither is indexed.
\code{logout\_view} still accepts GET. \code{UserCreationForm} is imported and unused,
\code{is\_admin} is defined twice, and \code{django.contrib.admin} is routed at
\code{/admin/} with nothing registered in it.
\end{honest}

\subsection{Why a green suite missed all of this}

\begin{honest}{The tests assert on page text, and on the easy branch}
Every assertion is a substring match against the whole HTML response, chrome included.
The project already has a scar from this: \code{assertNotContains(response, 'Max')} once
failed on a page with no donation for the pet Max, because the filter form's ``Max
Amount'' label contained the substring. It was resolved by renaming the label to ``Amount
To''. That fix is backwards, and worth admitting: the user interface was changed to suit
the test's matching strategy, and the test is still fragile.

Beyond fragility, three tests assert less than their names promise.
\code{test\_filter\_by\_payment\_method} checks that the cash donor appears but never
that the online donor is absent, so a filter that ignored its parameter entirely would
pass. \code{test\_amount\_validation} checks only that no row was created, so a view that
crashed with a 500 would pass. And \code{test\_csv\_export\_...} exercises the export
with no filter applied, which is the branch that works.

That last one is the whole lesson. Eleven green tests bought real confidence about
authentication and CRUD, and no confidence at all about the two features most likely to
put a wrong number, or somebody's phone number, in front of a user.
\end{honest}

\section{What is genuinely good}

Said plainly, because it is not padding. The custom user model was done on day one, which
is the one Django decision that is nearly impossible to walk back. Money is a
\code{Decimal}. The amount filters use \code{if amount\_min is not None} rather than a
truthiness test, a deliberate guard against falsy zero showing the author was thinking
about edge cases. The self-deletion guard is enforced in the view, not merely hidden in
the template. The top-pets panel once aggregated over the unfiltered table and
contradicted the total beside it; noticing and fixing that is exactly the class of bug
most people ship.

And the \code{TESTING = 'test' in sys.argv} flag is a real debugging win. Working out why
every test suddenly 301-redirects to https requires knowing that \code{manage.py test}
does \emph{not} set \code{DEBUG=True}, so a block guarded by \code{if not DEBUG} applies
during tests and \code{SECURE\_SSL\_REDIRECT} silently breaks the entire suite. That is a
non-obvious chain, and the README reconstructs it correctly.

The README volunteers several of its own weaknesses, including the missing foreign key,
the session export and the missing indexes. That is rarer than it should be and worth
defending as deliberate practice, which is exactly why the one place it gets a detail
backwards (the currency claim) is worth correcting before someone else finds it.

\section{If it continued}

\begin{enumerate}
  \item Fix the empty-filter export. Live data-exposure bug, one-line change, and add the
    test that catches it.
  \item Group totals by currency, or constrain the field and then group. Stop printing a
    number that is not true.
  \item Untrack \file{db.sqlite3}.
  \item Validate admin passwords; stop making every admin a superuser.
  \item Add \code{entered\_by} to \code{Donation}; stop hard-deleting.
  \item Re-derive the export from the query string and delete the session handoff, which
    retires two findings at once.
  \item Make \code{SECRET\_KEY} mandatory in production; index \code{donation\_date} and
    \code{pet\_name}; split settings into base, dev and prod.
\end{enumerate}

\begin{keyidea}{The one-sentence summary}
A small, conventional, competently built Django CRUD app that gets several non-obvious
Django decisions right and carries two real bugs in the two places nobody tested: an
empty filter exports everything, and mixed currencies are added together and labelled
PKR.
\end{keyidea}

\section*{Glossary}

\begin{description}[leftmargin=0pt, style=nextline, itemsep=3pt]
  \item[Django] A Python web framework: a pre-built skeleton handling the plumbing
    (listening for requests, matching URLs, talking to the database, rendering HTML) so
    you write only what is specific to your problem.
  \item[App] A Django folder grouping the models, views and templates for one area of
    responsibility. Not a separate program. One project contains many.
  \item[Model] A Python class describing one kind of record, which Django maps onto a
    database table.
  \item[View] A Python function that receives a request, does the work, and returns a
    response.
  \item[Template] An HTML file with placeholders filled in with data at render time.
    Django escapes values by default, which is what makes displaying donor-entered names
    safe.
  \item[ORM] Object-Relational Mapper. Lets you write
    \code{Donation.objects.filter(pet\_name='Max')} instead of SQL, and returns Python
    objects rather than rows.
  \item[QuerySet] What the ORM returns. Lazy: it records conditions and does not query
    the database until something reads the results.
  \item[Migration] A recorded, versioned instruction for changing the database structure,
    generated from the models and applied in order.
  \item[Middleware] Components every request passes through on the way in and every
    response on the way out, handling concerns no single view should repeat.
  \item[Session] Server-side storage tied to one browser by a cookie. The browser holds
    an opaque key; the data lives in a database table.
  \item[CSRF] Cross-Site Request Forgery: an attack where another site quietly submits a
    request to one you are logged into, riding your cookie. Django defends with a secret
    token in every form.
  \item[Foreign key] A column holding the identifier of a row in another table, expressing
    ``this donation was entered by that user''. The relationship this project lacks.
  \item[Decimal vs float] Decimals store digits exactly; floats store binary
    approximations and drift. Money uses Decimal.
  \item[UUID] A random 128-bit identifier used instead of a sequential integer, so ids
    cannot be enumerated or guessed.
  \item[Password hashing] Storing a one-way scrambled form of a password. Django uses
    PBKDF2-SHA256 at 600{,}000 iterations; the deliberate slowness makes guessing
    expensive.
  \item[Whitenoise] A library letting the Python app serve its own CSS and images
    competently, so no separate web server is needed.
  \item[collectstatic] A build step gathering static files into one directory; the
    manifest variant also hashes filenames for caching, and raises an error if the build
    step never ran.
  \item[Gunicorn] The production server that runs the Python application, named in the
    \file{Procfile}.
  \item[CRUD] Create, Read, Update, Delete.
\end{description}

\vfill
\begin{center}
\small\itshape
Every claim here was checked against the source, and every claim marked verified was
established by executing the code: the test suite, the form validation probes, the export
behaviour, the mixed-currency totals, and the contents of the committed database.
\end{center}

\end{document}
